% =============================================================================
% ARCHIVIST'S NOTE — Diegetic Frame of Archive 1,260
% =============================================================================
\cleardoublepage
\thispagestyle{empty}
\chapter*{Archivist's Note}
\addcontentsline{toc}{chapter}{Archivist's Note}
\markboth{Archivist's Note}{}

\begin{center}
\itshape\small
This envelope contains the complete transcription of the diary of Patient 1,260.\\
Reading the original material caused psychosomatic alterations in the examiners.\\[0.5em]
\upshape\scriptsize --- Fountain-pen annotation on forensic envelope no. 1,260/1979
\end{center}

\vspace{1.2em}

\begingroup

The volume you hold in your hands was not written to be published. It is the forensic reconstruction of the diary of \textbf{Joaquim Antônio Correia}, Patient 1,260 of the Barbacena Hospital-Colony.

A survivor of the 1915 drought and of the concentration camps for refugees in Ceará, Joaquim was committed to the Colony in 1917 and remained there for 62 years, until his death in 1979. Throughout that imprisonment, he produced 847 pages of fragmentary notes. The original notebook was found tied to his chest with strips of cloth when his body was exhumed from the mass grave --- 62 days after burial.

\vspace{0.4em}
\noindent According to the gravedigger's forensic account, the paper gave off a sweet odor of dampness and earth, and the ink seemed never to dry. The coroner's report registered the find as a custody object and attached it to the medical record.

\vspace{0.6em}
This edition brings together 84 fragments organized in five dimensions: the memories of childhood in the Sertão; the isolation and the tortures of the psychiatric hospital; the examination of the clinical reports; the forensic investigation conducted by Dr. Lúcia Mendes in 2026; and the record of the reader's own contamination.

\vspace{0.6em}
This is not a simple horror story. The pain documented in these pages is not fiction: it crossed decades, consumed doctors and examiners, and seeks, in the act of reading, a new host.

\vspace{1.2em}
\noindent{\itshape
It is recommended that you do not read this book in the dark.\\[0.35em]
It is recommended that you do not read this book in absolute solitude.\\[0.35em]
If the smell of wet earth arises during the reading, close the volume immediately.}

\vspace{1.2em}
\noindent The historical wound upon which this narrative rests is real. The Barbacena Hospital-Colony (1903--1994) interned more than 60,000 people without their consent; 1,853 bodies were exhumed from its mass graves. Joaquim is a character, but the pain he carries --- hunger, displacement, confinement, psychiatric torture --- belonged to thousands of anonymous Brazilians.

\vspace{0.6em}
\noindent The horror evoked here must not be consumed as entertainment. The manipulation of the reader is an ethical and formal choice --- not to deceive, but to make the body feel what the mind knows: that the trench is not a metaphor, and that the documented suffering exists to be witnessed, not aestheticized.

\vspace{0.6em}
\noindent{\itshape You will be manipulated. You will be contaminated. But in the end, you will know why.}

\vspace{1.4em}
\hfill{\scshape --- The Archivist}\hspace{1em}

\endgroup
\clearpage
